\section{Pivot: Visual Prompting iterativo en la navegación robótica}

\subsection{Planteamiento del problema y método}

Pivot resuelve el problema de la navegación visual robótica: el robot necesita navegar y desplazarse en el entorno según una instrucción (por ejemplo, «encuentra la banca de madera»).

\begin{importantbox}{La optimización iterativa de Pivot}
Pivot utiliza una estrategia de Visual Prompting iterativa:
\begin{enumerate}
    \item Muestrea un conjunto de direcciones candidatas a partir del espacio de acciones
    \item Mapea las acciones candidatas al espacio de la imagen (marca flechas de dirección o números sobre la imagen)
    \item Hace que el VLM infiera, con base en la imagen marcada, las mejores direcciones candidatas
    \item Vuelve a muestrear dentro del subespacio seleccionado y repite el proceso anterior
    \item Tras varias iteraciones, produce la acción final
\end{enumerate}
\end{importantbox}

Las entradas centrales del algoritmo:
\begin{itemize}
    \item Imagen RGB original
    \item Instrucción de navegación
    \item Definición del espacio de acciones
    \item Número máximo de iteraciones $T$
    \item Cantidad de muestras por ronda $K$
\end{itemize}

En cada iteración, se muestrean $K$ acciones a partir de una política de acción probabilizada $\pi$; tras mapearlas al espacio de la imagen, se hace que el VLM elija los mejores candidatos, y luego se ajusta una nueva distribución de acciones; el ciclo continúa hasta alcanzar el número máximo de iteraciones. En esencia, es un proceso de optimización de muestreo, selección y refinamiento.

\subsection{El valor del Visual Prompting en la navegación}

En las tareas de navegación, el valor del Visual Prompting radica en convertir la selección abstracta de acciones (avanzar, girar 30 grados a la izquierda, girar 15 grados a la derecha, etc.) en una selección visual directa: se marcan sobre la imagen los números de las direcciones candidatas, y se hace que el VLM elija la dirección más razonable con base en la información visual.

\begin{warningbox}{La particularidad del escenario de navegación}
A diferencia de la localización de objetos, la navegación no requiere elegir «qué objeto de la imagen», sino decidir «hacia qué dirección moverse». Aquí, la función del Visual Prompting es discretizar el espacio continuo de direcciones en un número finito de opciones marcadas, lo que reduce la dificultad de decisión del VLM.
\end{warningbox}

\subsection{Resumen de la sección}

Pivot muestra la aplicación del Visual Prompting en la navegación robótica: mediante un flujo iterativo de muestreo, marcado, inferencia del VLM y refinamiento, logra decisiones de navegación basadas en información visual, y es un trabajo representativo de la extensión del Visual Prompting de la localización estática a la toma de decisiones dinámica.

